% ============================================================
%  CHAPTER 5: PROPOSED WORK
%  Automated Validation of Static Vulnerability Candidates
%  Through Directed Symbolic Execution
% ============================================================
\chapter{Proposed Work: Automated Validation of Static Vulnerability Candidates Through Directed Symbolic Execution}
\label{ch:proposed}

\section{Introduction}
\label{sec:proposed-intro}

\mango (Chapter~\ref{ch:mango}) can analyze every binary in a firmware image and produce thousands of taint-style vulnerability alerts, each with Rich Expression evidence describing how attacker input transforms from source to sink.
On the \karonte dataset alone, \mango generates over 2,000 alerts across 49 firmware images.
However, none of these alerts come with proof: \rda reports data flows but cannot determine whether the conditional checks along those paths are satisfiable under any concrete input.
Triaging these candidates manually is infeasible at scale.

Symbolic execution~\cite{king1976symbolic} can automate this validation by producing concrete triggering inputs or proving infeasibility, but whole-program symbolic execution on firmware binaries is intractable due to path explosion~\cite{krishnamoorthy2010tackling} and environmental preconditions (NVRAM flags, proprietary HTTP APIs, cross-binary state) that generic executors cannot satisfy.

I observe that \mango's static analysis has already solved the hardest part: identifying \emph{which} paths to examine.
Each closure encodes a complete function-level trace from source to sink.
Rather than exploring the binary exhaustively, a symbolic executor can replay only these traces, transforming validation from an open-ended search into a bounded verification task.

To this end, I propose a directed symbolic execution system that takes \mango closures as input and produces a verdict (reachable, unreachable, or constant sink) with concrete proof-of-concept inputs for confirmed paths.
The system operates on angr's AIL, a decompiler-level IR that produces tighter constraints than VEX, and combines static and LLM-directed precondition recovery to handle firmware-specific guards.
I plan to evaluate this system across multiple firmware images, vendors, and architectures.

\section{Approach}
\label{sec:proposed-approach}

The proposed system takes \mango closures and a firmware binary as input.
Rather than exploring the binary exhaustively, the executor targets only the paths that static analysis already identified, avoiding the path explosion that makes conventional symbolic execution intractable.
For each closure, the system produces a verdict (reachable, unreachable, or constant sink) and, when reachable, a concrete proof-of-concept input.

The approach operates on angr's AIL (angr Intermediate Language), a higher-level IR that carries recovered type information and structured control flow from decompilation.
Operating at this level produces tighter constraints and fewer infeasible paths than VEX-based execution.
To handle firmware-specific guards, the system combines static precondition extraction (identifying NVRAM keys, HTTP parameters, and guard patterns from the decompilation) with LLM-directed precondition recovery for cases where static extraction is insufficient.
A hybrid concretization policy allows guard checks to resolve deterministically while preserving provenance annotations that track which attacker-controlled sources influence the sink argument.

For vulnerabilities that span multiple binaries through NVRAM, the system traces SET and GET locations across the firmware image using \mango's environment dictionary, constructing multi-stage exploit chains.

\section{Preliminary Evidence}
\label{sec:proposed-preliminary}

A prototype has been evaluated on a NETGEAR R7000 firmware image containing 161 closures across 33 binaries.
The prototype confirms 51 reachable command injection paths (31.7\%), identifies 50 closures with constant sink arguments (31.1\%), and produces concrete proof-of-concept exploits with canary validation for every confirmed path.
Median execution time per closure is 20.4 seconds.
These results demonstrate feasibility but are limited to a single firmware image from one vendor.

\section{Planned Evaluation}
\label{sec:proposed-evaluation}

The proposed evaluation will expand the prototype across multiple firmware images, vendors, and architectures to assess generalization.
I plan to conduct the following experiments:

\begin{enumerate}[nosep]
  \item \textbf{Multi-firmware effectiveness.} Evaluate on closure sets from additional firmware images spanning multiple vendors and architectures (ARM and MIPS). Report confirmation rates and compare against the R7000 baseline.
  \item \textbf{Comparison with existing approaches.} Compare confirmation rate and false-positive elimination against angr's VEX-based symbolic execution and Arbiter's static feasibility filtering on the same closure sets.
  \item \textbf{End-to-end exploit confirmation.} Execute generated PoC scripts against emulated firmware to validate that symbolic execution results correspond to exploitable vulnerabilities.
\end{enumerate}

\section{Timeline}
\label{sec:proposed-timeline}

\begin{center}
\begin{tabular}{ll}
\toprule
\textbf{Year} & \textbf{Milestone} \\
\midrule
2020 & Began Ph.D. at Arizona State University (SEFCOM) \\
2022 & Passed comprehensive exam \\
2024 & Published Operation Mango (USENIX Security 2024) \\
2025 & Artiphishell at DEF CON (DARPA AIxCC); prospectus defense \\
2025--2026 & Multi-firmware evaluation and baseline comparisons \\
2026 & Paper submission; dissertation writing and defense \\
\bottomrule
\end{tabular}
\end{center}
